\documentclass[12pt,letterpaper]{article}
\usepackage{amsmath,amssymb,amsthm}
\begin{document}
\begin{center}
{\Large\textbf{Fourier Series, Briefly}}\\[6pt]
\small Lecture notes
\end{center}
\bigskip

Let $f$ be $2\pi$-periodic and integrable. Its Fourier coefficients are
\begin{align}
a_k&=\frac1\pi\int_{-\pi}^{\pi}f(x)\cos kx\,dx,\\
b_k&=\frac1\pi\int_{-\pi}^{\pi}f(x)\sin kx\,dx,
\end{align}
and the series is $\frac{a_0}2+\sum_{k\ge1}(a_k\cos kx+b_k\sin kx)$.

\medskip
\noindent\textbf{Example.} For $f(x)=x$ on $(-\pi,\pi)$ we get $a_k=0$ and $b_k=2(-1)^{k+1}/k$. Evaluating at $x=\pi/2$ yields the Leibniz series
\[
\frac\pi4=1-\frac13+\frac15-\frac17+\cdots.
\]

\vspace{1ex}
\noindent\textbf{Parseval.} For square-integrable $f$,
\begin{equation*}
\frac1\pi\int_{-\pi}^{\pi}|f|^2\,dx=\frac{a_0^2}2+\sum_{k=1}^\infty(a_k^2+b_k^2).
\end{equation*}
Applied to $f(x)=x$ this gives $\sum_{k\ge1}k^{-2}=\pi^2/6$.
\hrule
\vspace{2pt}
\flushright{\small End of notes.}
\end{document}
